\documentclass[10pt,twocolumn]{article}
\usepackage[a4paper, left=1.5cm, right=1.5cm, top=2cm, bottom=3cm]{geometry}
\usepackage[T1]{fontenc}
\usepackage[utf8]{inputenc}
\usepackage[italian]{babel}
\usepackage{amsmath}
\usepackage{titling}
\usepackage{amssymb}
\usepackage{caption}
\usepackage{graphicx}
\usepackage{float}
\usepackage{relsize}
\usepackage{amsmath}
\usepackage{titlesec}
\usepackage{ragged2e}
\usepackage{circuitikz}
\usepackage{booktabs}
\usepackage{longtable}
\usepackage{makecell}
\usepackage{pdflscape}
\usepackage{enumitem}
\usepackage{tikz}
\usepackage{physics}
\usepackage{xcolor}
\usepackage[most]{tcolorbox}
\usepackage{tikz-3dplot}
\usepackage{tikz}
\usepackage{ragged2e}
\usepackage{siunitx}
% \usepackage{booktabs}
\usepackage[colorlinks=true, linkcolor=black]{hyperref}  %per rendere l'indice genrale "interattivo"
\usepackage{enumitem}  %distanza degli itemize
\setlist[itemize]{itemsep=4pt, parsep=1pt}
\newtcolorbox{nota}{
  blanker,
  before skip=1em,
  after skip=1em,
  left=1em,
  borderline west={1pt}{0pt}{black},
  fontupper=\itshape,
  before upper={\noindent\textbf{Nota}:\quad}
}

\setlength{\intextsep}{8pt} % per regolare la distanza grafico tabelle
\setlength{\columnsep}{60pt}
\setlength{\tabcolsep}{4pt}

\begin{document}

    \justifying 

      \title{\textbf{Verifica delle legge delle lenti sottili}}
      \author{Tedeschini Michele \hspace{0.7cm} Stroili Emanuele \hspace{0.7cm} Mirolo Manuele \hspace{0.7cm} Brusini Alessio} 
      \date{16 Ottobre 2026}
      \maketitle
      \newgeometry{left=3cm, right=3cm, top=4cm, bottom=4cm}
      \onecolumn
      \tableofcontents
    \vspace{3cm}
      \begin{abstract}
        \centering
        \large
        L'esperimento consiste nel verificare l'equazione di Sellmeier per l'indice di rifrazione di un prisma. 
          
      \end{abstract}

      \newpage
    \restoregeometry
    \twocolumn



    \section{Legge fisica da verificare}
    \section{Apparato sperimentale}
    \section{Grafici}
    \section{Conclusioni}
    \end{document}